\section{Scope and status}
The motivating question is whether the information in a complete table of SAT
answers can determine the least computational effort needed to answer a new
SAT query. We grant access to the complete finite table and do not charge for
discovering a representation. This convention makes the resulting model
nonuniform: different input lengths may receive unrelated descriptions.
The object of study is the feasible relationship between description
length and query time.

The initial version is an expository foundation and a research protocol.
Its proved statements are standard consequences of computability, advice
complexity, counting, and residual-function arguments. No originality is
claimed for them. Proofs expose model assumptions and quantifiers.
The proposed lower bound for unrestricted local descriptions is an open
target, not a result of the paper.

Three quantities remain separate: the length of a stored description,
the work space of an evaluator, and its running time. A short program can
describe a long computation, while a computation can reuse its work space.
Exact deterministic bit recovery is a particular notion of prediction;
it is not a claim about probabilistic prediction, average log-loss, or
Shannon entropy. No unrestricted equivalence between those statistical
quantities and worst-case SAT running time is assumed.

Every substantive statement has a stable identifier, recorded in the
repository claim register. The register distinguishes expository proofs,
literature results with hypotheses, open targets, and finite observations.
Independent human review and proof-assistant verification are not claimed.
